%% kapa: exam=1 round=27 year=2016 period=1 form=A subject=civil_law subject_ko=민법 seq=35 number=35
%% source: Q-Net 감정평가사 기출문제 (cst00309) artlSeq=5169402 file=extracted/1008/5169402/제27회 감정평가사 1차 문제/제27회 감정평가사 1교시문제(A형).pdf
%% conversion: kordoc 4.24.51 -> tools/gen1.py (auto); review=proofread-sonnet; flags=none
\begin{question}{35}
유치권에 관한 설명으로 옳지 않은 것은? (다툼이 있으면 판례에 따름)
\begin{choices}
  \item 임차인의 임차보증금반환청구권은 임차건물에 관하여 생긴 채권이라 할 수 없다.
  \item 점유를 침탈 당한 유치권자가 점유회수의 소를 제기하면 유치권을 보유하는 것으로 간주된다.
  \item 유치권의 발생을 배제하는 특약은 유효하다.
  \item 피담보채권이 변제기에 이르지 아니하면 유치권을 행사할 수 없다.
  \item 유치권자는 유치물의 과실을 수취하여 다른 채권보다 우선하여 그 채권의 변제에 충당할 수 있다.
\end{choices}
\end{question}
\answer{2}
